\chapter*{Tóm tắt}
%\markboth{TÓM TẮT}

Đồ án môn học này tập trung vào việc nghiên cứu và cải thiện khả năng tự phục hồi (self-healing) cho hệ thống Kubernetes sử dụng Học tăng cường (Reinforcement Learning - RL), nhằm thay thế cho các giải pháp cấu hình tĩnh truyền thống. Chúng em mô hình hóa bài toán tự phục hồi dưới dạng tiến trình quyết định Markov (MDP) với không gian trạng thái từ các chỉ số đo lường hệ thống (metrics), không gian hành động can thiệp (khởi động lại pod, co giãn tài nguyên) và hàm thưởng cân bằng giữa tài nguyên và chất lượng dịch vụ (SLA). 

Ba thuật toán học tăng cường tiêu biểu gồm DQN, A2C và PPO được cài đặt và huấn luyện trên môi trường giả lập. Kết quả đánh giá cho thấy tác nhân PPO đạt tính hội tụ ổn định và tỷ lệ tự phục hồi thành công cao nhất so với các thuật toán còn lại và phương pháp baseline. Khi triển khai thử nghiệm trên cụm Kubernetes thực tế chạy nền tảng AWS, mô hình đề xuất đã giúp cải thiện đáng kể các chỉ số chất lượng dịch vụ (SLA) như thời gian phản hồi và tỷ lệ lỗi, trong khi mức tiêu thụ tài nguyên của cụm vẫn được duy trì ở mức tối ưu.
